\section{The microscopic boundary and relative factorization}\label{supp:sec:boundary}
The exact border cylinder and microscopic non-vacancy are different
objects. The first has probability $2^{-\pi(L)}$ without an error term;
the second also allows an interior start. This appendix proves that the
entire interior first moment is exponentially smaller than the border
probability. A union bound is then sufficient; no microscopic overlap
expansion is needed.

\subsection{Valuation matrices and the transition range}
Put $B=L+1$, $V_x=\{x-1,\ldots,x+L-1\}$, and let $M_x$ have rows
$v(n)=(v_p(n)\bmod2)_p$ for $n\in V_x$.
The start matrix is $A_x=D_xM_x$, where $D_x$ is the edge-difference
matrix of a tree on $V_x$. Its kernel consists exactly of the constant
vectors. Consequently
\begin{equation}\label{supp:eq:boundary-rank}
 \rank A_x\ge\rank M_x-1,\qquad
 \PP(J_{x,L}=1)\le2^{-\rank A_x}.
\end{equation}
For $2\le x\le L-1$ a start is impossible: both $2x$ and $2x-2$
lie inside the putative constant run, and complete multiplicativity
would force $f(2)=1$ and $f(2)=-1$ simultaneously.

\begin{lemma}[Transition identity minor]\label{supp:lem:transition}
Uniformly for $x\in\{B-1,B,B+1\}$,
$\rank A_x\ge(2-o(1))\pi(B)$.
\end{lemma}
\begin{proof}
Write $a=x-1\in\{B-2,B-1,B\}$ and $Y_0=2\sqrt B$.
For each prime $Y_0<p\le B$ choose
$m_p=p\lceil a/p\rceil\in[a,a+B-1]$.
The cofactor satisfies
$\lceil a/p\rceil\le B/p+1<\sqrt B/2+1<Y_0<p$ for large $B$.
Thus $m_p$ has valuation one at $p$ and no other prime factor above
$Y_0$. Each prime $B<q\le a+B-1$ occurs as its own vertex.
On these rows and columns the valuation matrix has an identity minor
of size $\pi(a+B-1)-\pi(Y_0)=(2-o(1))\pi(B)$.
Apply \eqref{supp:eq:boundary-rank}.
\end{proof}

\subsection{The uniform prime-divisor input}
The precise literature input is
\begin{equation}\label{supp:eq:LS-uniform}
 \forall\varepsilon>0\ \exists B_0(\varepsilon)\ \forall B\ge B_0(\varepsilon)
 \ \forall n>B:\qquad
 \omega\left(\prod_{j=0}^{B-1}(n+j)\right)\ge(2-\varepsilon)\pi(B).
\end{equation}
It is equation~(14) of Laishram--Shorey \cite{LaishramShorey2004}.
Its uniformity is important. For $n>17B/12$, their Theorem~2 gives
at least $\pi(2B)$ prime divisors, apart from a fixed exceptional pair.
For $B<n\le17B/12$, every prime at most $B$ divides the block product,
and each prime in $[n,n+B-1]$ is an additional divisor. Hence the number
is at least
\[
 \pi(B)+\pi(n+B-1)-\pi(n-1)=(2-o(1))\pi(B),
\]
uniformly in that bounded ratio, by the prime number theorem.
This verifies the quantifier order in \eqref{supp:eq:LS-uniform}.

Suppose now that $B+2\le x\le B^2+2$. Every vertex is at most
$B^2+B$. Equation~\eqref{supp:eq:LS-uniform} supplies
$t_x\ge(1-o(1))\pi(B)$ primes $q>B$ in the block.
Each has valuation one, since $q^2\ge(B+1)^2>B^2+B$; it occurs at
one vertex only, since $q>B-1$; and two such primes cannot occupy the
same vertex, since their product exceeds $B^2+B$.
Their columns span a coordinate subspace $U_x$ of dimension $t_x$.
Call its rows the private rows.

\subsection{A global count of even valuations}
Fix an integer $K\ge2$ and use the medium primes $B/K<p\le B$.
Let $E_x$ be their total odd-valuation incidence mass in $V_x$.
We claim
\begin{equation}\label{supp:eq:medium-mass}
 E_x\ge(H_K-1-o(1))\pi(B),\qquad H_K=\sum_{j=1}^K\frac1j.
\end{equation}
The block contains at least $\lfloor B/p\rfloor$ multiples of $p$.
An incidence of positive even valuation is among those with
$n=ap^2\le B^2+B$. Thus $a\le K^2+1$ for sufficiently large $B$.
For fixed $a$, successive distinct primes $p,q>B/K$ give values
$ap^2,aq^2$ separated by more than $2aB/K$; only
$1+K/(2a)$ such values can lie in an interval of length $B$.
Summing over the fixed number of possible $a$ shows that the total
number of square divisibilities is $O_K(1)$, uniformly in $x$.
Subtracting all of them is a permissible overcount of even valuations.

The main incidence count can be evaluated without rounding every summand:
\begin{align*}
 \sum_{B/K<p\le B}\left\lfloor\frac Bp\right\rfloor
 &=\sum_{j=1}^{K-1}\bigl(\pi(B/j)-\pi(B/K)\bigr)\\
 &=(H_{K-1}-(K-1)/K+o(1))\pi(B)\\
 &=(H_K-1+o(1))\pi(B).
\end{align*}
This proves \eqref{supp:eq:medium-mass}. The improvement over a
prime-by-prime loss is essential: subtracting one for each prime would
lose a positive proportion of $\pi(B)$.

\subsection{Rank of a simple incidence matrix}
\begin{lemma}[Simple incidence rank]\label{supp:lem:simple-incidence}
Let $H$ be a finite binary matrix. Suppose every row has weight at most
two, every column has weight at most $K$, and no pair of columns occurs
as the support of two weight-two rows. If $E$ is the number of entries
equal to one, then
\[
 \rank H\ge\frac{E}{K+1}.
\]
\end{lemma}
\begin{proof}
Remove zero columns and view a weight-two row as an edge between its
columns, and a weight-one row as a pin. In a connected component with
$r$ active columns, an unpinned component has rank $r-1$ and a pinned
component has rank $r$. The first assertion follows because the edge
vectors of a connected graph span the even-parity subspace; one pin
then spans its missing direction.

A pinned component has incidence mass at most $Kr$, which is at most
$(K+1)$ times its rank. An unpinned component has $r\ge2$; simplicity
and the degree bound give incidence mass at most
$\min\{r(r-1),Kr\}$. If $r\le K+1$ use the first bound, and if
$r\ge K+1$ use the second. In either case the mass is at most
$(K+1)(r-1)$. Sum over the components.
\end{proof}

Delete the private rows from the medium-prime columns of $M_x$, and
call the resulting matrix $H_x$. For fixed $K$ and sufficiently large
$B$, a private row carries at most one medium prime: otherwise a vertex
would exceed $B(B/K)^2>B^2+B$. Thus deletion removes at most $t_x$
incidences. A remaining row carries at most two medium primes, since
$(B/K)^3>B^2+B$. A column has weight at most $K$.
Finally two weight-two rows with the same support $p,q$ would be two
multiples of $pq>(B/K)^2>B$ inside one block, which is impossible.
All hypotheses of \cref{supp:lem:simple-incidence} are therefore met.
Modulo $U_x$, a medium column is represented by its restriction to the
remaining rows, so
\begin{align}\label{supp:eq:surplus-K}
 \rank A_x
 &\ge t_x+\frac{(E_x-t_x)_+}{K+1}-1\notag\\
 &\ge \left(1+\frac{H_K-2}{K+1}-o(1)\right)\pi(B).
\end{align}
For the last step, if $E_x\le t_x$ use $t_x\ge E_x$; otherwise use
$t_x\ge(1-o(1))\pi(B)$ and \eqref{supp:eq:medium-mass}.
The parameter can be chosen exactly within this family. Set
$\delta_K=(H_K-2)/(K+1)$. Then
\[
 \delta_{K+1}-\delta_K
   =\frac{3-H_K}{(K+1)(K+2)}.
\]
Since $H_{10}<3<H_{11}$ and $H_K$ is strictly increasing, the unique
maximizer over integers $K\ge2$ is $K=11$. Consequently we use
\[
 H_{11}=\frac{83711}{27720},\qquad
 \delta_*:=\frac{H_{11}-2}{12}
          =\frac{28271}{332640}>\frac1{12}.
\]
This optimizes only the displayed fixed-cutoff family of bounds,
not the rank surplus among all possible methods.

The local product comparisons hold for $B\ge11^3+1=1332$:
$(B/11)^3>B^2+B$, $B^3/121>B^2+B$ and $(B/11)^2>B$.
In the square-divisibility count one may then use $a\le122$ and at most
seven values per $a$, giving the safe bound $854$. These elementary
checks do \emph{not} make $1332$ a threshold for the complete asymptotic
theorem: the prime-divisor and prime number theorem thresholds still
depend on the desired error margin.

\subsection{Post-quadratic starts and summation}
\begin{lemma}[Post-quadratic defect maximum]\label{supp:lem:pointwise-maximum}
There is an absolute constant $\theta_0$ such that, for all sufficiently
large $B$ and every $x>B^2+2$, the number $m_x$ of $B$-defective
vertices in $V_x$ is at most $B-g_L$, where
\[
 g_L=\frac B{\log B}
       (\log\log B-\log\log\log B-\theta_0),
 \qquad g_L/\pi(B)\longrightarrow\infty.
\]
In particular $\PP(J_{x,L}=1)\le2^{-g_L}$.
\end{lemma}
\begin{proof}
We use the Balasubramanian--Shorey exclusion
\cite{BalasubramanianShorey1993} in the form recorded by Shorey
\cite[pp.~55--56, equation (15) with exponent $2$ and the subsequent
square-case statement]{Shorey1994BakerApplications}.
It has an absolute constant $\theta_0$ and, for every fixed $\varepsilon>0$,
a threshold $k_0(\varepsilon)$ such that, for $k\ge k_0(\varepsilon)$,
\[
 \prod_{i=1}^t(z+d_i')=by^2,\quad
 0\le d_1'<\cdots<d_t'<k,\quad b>0,\quad P^+(b)\le k,
 \quad z>e^{1-\theta_0+\varepsilon}\frac{k\log k}{\log\log k}
\]
forces $t<\mu_k(\theta_0)$, where
\[
 \mu_k(\theta)=k\left(1-\frac{\log\log k}{\log k}
       +\frac{\log\log\log k}{\log k}+\frac\theta{\log k}\right).
\]
The smoothness condition on $b$ is part of (15); it is not a size
bound on $b$. All thresholds are uniform in $z,b,y$ and the shifts.
To obtain the form used here, first fix $\varepsilon=1$. Increase the
absolute threshold until
\[
 k^2>e^{2-\theta_0}\frac{k\log k}{\log\log k}.
\]
For $m>k^2$ and $1\le d_1<\cdots<d_t\le k$, set
$z=m+1$ and $d_i'=d_i-1$. These substitutions preserve every factor
$m+d_i=z+d_i'$, put all shifts in $[0,k)$, and satisfy the required
height bound. Hence the product equation forces $t<\mu_k(\theta_0)$.
Write $V_x=\{m+1,\ldots,m+B\}$ with $m=x-2>B^2$.
The product of its $B$-defective vertices is a square times a
$B$-smooth integer. There is no extra upper bound on that smooth
coefficient in the imported statement. Its contrapositive gives
$m_x<\mu_B(\theta_0)=B-g_L$ beyond the fixed absolute threshold.
The affine defect bound is
$\PP(J_{x,L}=1)\le2^{-L}2^{\max(m_x-1,0)}$, which is at most
$2^{-g_L}$ for large $B$.
\end{proof}

Partition the starts $2\le x\le2L^2$ into the forbidden early range,
the three transition starts, $B+2\le x\le B^2+2$, and the remaining
post-quadratic starts. Their total probabilities are respectively
\[
 0,\qquad 3\,2^{-(2-o(1))\pi(B)},\qquad
 O(B^2)2^{-(1+\delta_*-o(1))\pi(B)},\qquad
 O(L^2)2^{-g_L}.
\]
Since $\log B=o(\pi(B))$ and $\pi(B)=\pi(L)+O(1)$, their sum is
at most $2^{-(1+\delta_*-o(1))\pi(L)}$.
Adding the exact border event proves
\[
 q_L=2^{-\pi(L)}\bigl(1+O_{\delta_0}(2^{-\delta_0\pi(L)})\bigr)
 \quad(0<\delta_0<28271/332640),
\]
including the displayed choice $\delta_0=1/12$.
The same interior sum divided by $q_L$ bounds the conditional
probability that the exact border is not the unique microscopic start.

The cutoff-invariance and mesoscopic estimates use one further
observation. At heights above $2L^2$, two $B$-defective vertices at
distance at most $L$ have distinct squarefree $B$-smooth parts; equal
parts would place consecutive squares too far apart. The number of
choices of these parts is $\exp(O(L/\log L))$, and the uniform Pell
bound gives the same type of bound on the number of corresponding windows in the entire admissible population
up to twice the ambient height. Multiplication by $2^{-g_L}$ remains
$o(2^{-\pi(L)})$, without an exterior height-slice factor. Starts with at most one defect
cost only their number times $2^{-L}$. This proves the bounds in the
article with cutoff $K\le2^L$, using $2^L$ as an artificial upper
ambient scale; it does not require a dyadic lower endpoint comparable
to that artificial scale.

\subsection{The prime clock}
Section~\ref{sec:prime-clock} of the article proves the exact prime-clock
law and explains why integer overshoots need not be tight. We record here
the conventions needed for the conditional comparisons. Let
$J_\partial=\min\{j:f(p_j)=-1\}$. On $C_L=1$ and the full-measure
event $J_\partial<\infty$, put $G_L=J_\partial-\pi(L)-1$;
assign $G_L=0$ elsewhere. Then
\[
 \PP(G_L=k\mid C_L=1)=2^{-k-1},\qquad
 \PP(\ell_\partial-L\ge h\mid C_L=1)
             =2^{-[\pi(L+h)-\pi(L)]}.
\]
Since $\{C_L=1\}\subseteq\{B_L^{\rm mic}>0\}$, the conditional
law of $G_L$ given microscopic non-vacancy is a mixture whose
$C_L=1$ component is exactly geometric. Its distance from that law is
at most $\PP(C_L=0\mid B_L^{\rm mic}>0)=O(2^{-\delta_0\pi(L)})$.
The next subsection constructs the finite-cylinder representative used
to carry this mark through the stable product comparison.

\subsection{Relative bulk estimates and affine conditioning}
\label{supp:sec:boundary-conditioning}
Let $\lambda=M2^{-L}\to0$, $L=O(\log M)$ and
$U=[\lceil M^\delta\rceil,M-L+1]\cap\Z$. Fix $\delta$ and an
upper bound for $L/\log M$. With $Y=Y_M^\sharp$, truncate the signed
exact marks at $E$ and put $Q=L+E+1$. The averaged conditional
process error is at most
\[
 C\left[p\mathcal M_B+p|D_Y(Q)|
       +p^2\bigl(MQ+E_Y(Q)+R_{2,\delta}(M,L)\bigr)\right].
\]
The local ranks are those in Appendix C. Summing signs and exact marks
before the separated-pair estimate leaves the relation sum at the base
length $L$; only the bad supports and shared-prime graph use $Q$.
After division by $Mp=\lambda$, the contributions are
\begin{center}\small
\begin{tabularx}{\textwidth}{@{}lX@{}}\toprule
Source & Normalized contribution\\\midrule
One-window defects & $M^{-1/2+o(1)}$\\
Bad maximal supports & $e^{-V_M+o(\nu_M)}$\\
Diagonal and local pairs & $pQ$\\
Shared-prime edges & $\lambda(e^{-V_M+o(\nu_M)}+M^{-1+o(1)})$\\
Base-length relations & $M^{-1/3+\varepsilon}(1+\lambda)$\\
Actual and target mark tails & $O(2^{-E-1})(1+M^{-1/2+o(1)})$\\\bottomrule
\end{tabularx}
\end{center}
For the last row, the exact tail event is a start at length $L+E+1$;
the masked first-moment estimate at that length gives the displayed
bound after averaging the conditional tail probability. Choose
$E=\lceil V_M/\log2\rceil$. Then $Q=O(\log M)$,
$pQ=O(\log M/M)$ for large $M$, and every row tends to zero.
This proves the signed, marked relative error in
\eqref{eq:relative-bulk-rate}, not merely its unmarked projection.

The microscopic vector is measurable at the same prime cutoff because
$2L^2+L<Y$ for large $M$. For a fixed $K$, define $G_{L,K}$ from the
first $K$ prime signs after $L$: count consecutive positive signs, capped
at $K$, and put the value zero off $C_L=1$.
This finite-cylinder variable agrees almost surely with $G_L\wedge K$;
this makes measurability on the finite prime sigma-field explicit,
independently of the convention on the null nonterminating event. Bertrand's postulate bounds the $K$th
prime after $L$ by $2^K L<Y$ eventually, for fixed $K$.
The stable lift retains $(\mathbf B_L,G_{L,K})$ jointly with the full
signed marked bulk field, at cost $\Delta_M=o(\lambda)$.

Write $\alpha=2^{-\pi(L)}$, $b=|U|p$, and let $u_M$ and $v_M$ denote
the interior microscopic first moment and the mesoscopic hit probability.
The source bounds give $u_M+v_M=o(\alpha+b)$.
The one-hit calculation in the proof of \cref{thm:two-clock}
gives sparse-target mass $e^{-b}[\alpha+(1-\alpha)b]$ and a change
of mixture weights at most $b$ relative to \eqref{eq:two-clock-target}.
The conditional bulk mark, sign and position retain their full laws;
only the boundary clock is truncated.
Simultaneous sources and multiple bulk points have probabilities at most
$\alpha b$ and $b^2/2$. Since the sparse-event mass is asymptotic to
$\alpha+b$, normalizing the source subprobability measures gives
\[
 \dTV\bigl(\Law(\Gamma_{M,L}^{[K]}\mid R_M\ge L),Q_{M,L}^{[K]}\bigr)
 \ll\frac{\Delta_M+u_M+v_M+\alpha b+b^2}{\alpha+b}+b.
\]
The additional $b$ is already absorbed by the last two numerator terms.
For the actual first-hit law the discarded boundary tail is at most
$2^{-K-1}$: its unnormalized mass is at most $\alpha2^{-K-1}$ and the
hit probability is at least $\alpha$. The target tail is also at most
$2^{-K-1}$. First let $M\to\infty$, then let $K\to\infty$.
No lower bound on either mixture weight, convergence of their ratio, or
exact finite-size independence is required. The grid-to-Lebesgue passage
is a separate weak limit.

Now let $A=\{GX=b\}$ be a compatible affine cylinder on the exposed
prime signs. Write $r=\rank G$, $I=r\log2$, and
$U_L=\Span\{e_p:p\le L\}$. If the stack with $C_L$ is compatible,
its rank is $r+\pi(L)-\kappa$, where
$\kappa=\dim(\operatorname{row}(G)\cap U_L)$.
Consequently
\[
 \PP(C_L=1\mid A)=2^{-(\pi(L)-\kappa)}
\]
exactly, whenever the exposed cutoff contains the border coordinates.
This finite identity needs no rare-intensity regime. If the stack is
incompatible, that probability is zero.
The parameter $\kappa$ counts consequences entirely supported on the
border coordinates, not input equations that merely mention them.

For the asymptotic conclusions assume $M,L\to\infty$,
$L\le A_0\log M$, $M2^{-L}\to0$, and eventually
$I\le V_M-c\nu_M$ for fixed $c>0$.
The matrices, their dimensions and right-hand sides may vary with $M$.
Multiply the quantitative normalized bulk ledger by $e^I$ before taking
limits. The cutoff term is at most $e^{-c\nu_M+o(\nu_M)}$; all polynomial
terms still vanish because $I=O(V_M)=o(\log M)$.
Thus the amplified bulk error is $o(\lambda)$; amplification of an
unspecified little-oh alone would not justify this conclusion.
Also $I=O(V_M)=o(\pi(L))$ and $\kappa\le r$. The conditioned
microscopic interior mass, relative to $2^{-\pi(L)+\kappa}$, is at
most
\[
 \exp(I)2^{-\delta_0\pi(L)-\kappa}=o(1).
\]
The mesoscopic first term is $\lambda M^{\delta-1}$ and remains
$o(\lambda)$ after multiplication by $e^I$; its other term is dominated
by $2^{-g_L}$ and remains negligible relative to the border mass.
The same one-hit calculation, normalized by $2^{-\pi(L)+\kappa}+b$,
therefore proves the affine-conditioned crossover with $\pi(L)$ replaced
by $\pi(L)-\kappa$. No assumption on future prime coordinates is used
for the rare tail, border concentration or location laws.

Finally let $q_1<\cdots<q_K$ be the first primes after $L$, and put
$Q_{L,K}=\Span\{e_{q_1},\ldots,e_{q_K}\}$ and
$H_L=\operatorname{row}(G)+U_L$. If the extended stack is compatible,
\[
 \PP(G_L\ge K\mid A,C_L=1)
                 =2^{-[K-\dim(H_L\cap Q_{L,K})]};
\]
otherwise this tail probability is zero. Thus the ordinary geometric
mark is preserved when, for every fixed $K$, the condition
$H_L\cap Q_{L,K}=\{0\}$ holds beyond a threshold that may depend on $K$.
This is separate from the parameter $\kappa$ controlling the border mass.
Fix $K$, use the finite-cylinder representative and joint stable lift,
and take the size limit before $K\to\infty$. No common threshold in $K$
and no neutrality for a growing truncation $K(M)$ is assumed.
The exact rank formula remains the alternative when neutrality fails;
no relative estimate for a disappearing mixture component is inferred.
